\scanpage{111}
\margin{86}
Utilizing (14) and assuming $n$ to be sufficiently large, we find from (21) and (22):
\[
\dim H^1(V,O_V[nE'])=\dim H^1(V,O_V(-\mathcal E'_\nu)[nE']),\qquad (23)
\]

We can assume that $m'\ge n(\mathcal E'_\nu)$ and so great that, at $n=m'$, (23) is fulfilled. Then we have an exact sequence
\diagram{\begin{tikzcd}
H^0(V,O_V[mE'])\arrow[r,"r^*"] & H^0(\mathcal E'_\nu,O_{\mathcal E'_\nu}[m'E'])\arrow[r] & 0
\end{tikzcd}}

Utilizing the properties of the homomorphism $r^*$(See p. 208) we obtain hence that the mapping of the variety $\mathcal E'_\nu$ into a projective space, corresponding to the group $\unsure{H^0(\mathcal E'_\nu,O_{\mathcal E'_\nu}[m'E'])\mathcal E'_\nu)}$, coincides with the restriction on $\mathcal E'_\nu$ of the mapping $\unsure{f_{m'}|f_m'}$. In particular, we obtain that $\unsure{f_{m'}|f_m'}$ maps $\mathcal E'_\nu$ in a one-to-one manner (that is, different points are transfered into different points).

Let $q'$ be so great that $r=2^{q'}m-m'>0$. Let us show that $f_{2^{q'}m}$ maps $\mathcal E'_\nu$ in a one-to-one manner. Let $P_1$ and $P_2$ be any two points on $\mathcal E'_\nu$. It is sufficient for us to show that there exists a cross-section of $H^0(V,O_V[2^{q'}mE'])$, which is equal to zero in $P_1$, but not in $P_2$. It follows from above that there
\mspage{110}
